\documentclass[11pt,a4paper]{report}

\usepackage[utf8]{inputenc}
\usepackage[T1]{fontenc}
\usepackage[english]{babel}
\usepackage[margin=2.5cm]{geometry}
\usepackage{titlesec}
\usepackage{enumitem}
\usepackage{xcolor}
\usepackage{tcolorbox}
\tcbuselibrary{skins,breakable}
\usepackage{fancyhdr}
\usepackage{tabularx}
\usepackage{longtable}
\usepackage{booktabs}
\usepackage{graphicx}
\usepackage{hyperref}
\usepackage{parskip}
\usepackage{lastpage}
\usepackage{listings}

\graphicspath{{../img/}}

% ---------- Colors ----------
\definecolor{primaryblue}{RGB}{20,50,90}
\definecolor{accentorange}{RGB}{214,110,40}
\definecolor{lightgray}{RGB}{242,242,242}
\definecolor{textgray}{RGB}{80,80,80}
\definecolor{codebg}{RGB}{247,247,247}

% ---------- Hyperref ----------
\hypersetup{
  colorlinks=true,
  linkcolor=primaryblue,
  urlcolor=primaryblue,
  citecolor=primaryblue,
  pdftitle={aws-iam-monitor v1.2.0 -- Full Documentation},
  pdfauthor={ZouariOmar},
  pdfsubject={aws-iam-monitor -- Full Technical Documentation},
  bookmarksnumbered=true,
}

% ---------- Code listings ----------
\lstset{
  basicstyle=\ttfamily\small,
  backgroundcolor=\color{codebg},
  frame=single,
  rulecolor=\color{textgray},
  breaklines=true,
  columns=fullflexible,
  showstringspaces=false,
}

% ---------- Titles ----------
\titleformat{\chapter}[display]
  {\normalfont\bfseries\color{primaryblue}}
  {\Large\filright\chaptertitlename~\thechapter}
  {6pt}
  {\Huge}
\titlespacing*{\chapter}{0pt}{0pt}{24pt}

\titleformat{\section}
  {\normalfont\Large\bfseries\color{primaryblue}}
  {\thesection}{0.6em}{}[{\color{primaryblue}\titlerule[0.8pt]}]
\titleformat{\subsection}
  {\normalfont\large\bfseries\color{primaryblue}}
  {\thesubsection}{0.6em}{}
\titleformat{\subsubsection}
  {\normalfont\bfseries\color{textgray}}
  {\thesubsubsection}{0.6em}{}
\titlespacing*{\section}{0pt}{16pt}{8pt}
\titlespacing*{\subsection}{0pt}{12pt}{6pt}

% ---------- Header / footer ----------
\pagestyle{fancy}
\fancyhf{}
\renewcommand{\headrulewidth}{0.4pt}
\renewcommand{\footrulewidth}{0.4pt}
\fancyhead[L]{\small\color{textgray}\nouppercase{\leftmark}}
\fancyhead[R]{\small\color{textgray} aws-iam-monitor v1.2.0}
\fancyfoot[C]{\small\color{textgray} Page \thepage\ / \pageref{LastPage}}

\fancypagestyle{plain}{
  \fancyhf{}
  \fancyfoot[C]{\small\color{textgray} Page \thepage\ / \pageref{LastPage}}
  \renewcommand{\headrulewidth}{0pt}
}

% ---------- Utility boxes ----------
\newtcolorbox{keybox}[1]{
  enhanced, breakable,
  colback=lightgray, colframe=primaryblue,
  boxrule=0.5pt, arc=2pt,
  left=8pt, right=8pt, top=6pt, bottom=6pt,
  title={#1}, coltitle=white, colbacktitle=primaryblue,
  fonttitle=\bfseries,
}

\newtcolorbox{notebox}[1][Note]{
  enhanced, breakable,
  colback=white, colframe=accentorange,
  boxrule=0.6pt, arc=2pt,
  left=8pt, right=8pt, top=5pt, bottom=5pt,
  title={#1}, coltitle=white, colbacktitle=accentorange,
  fonttitle=\bfseries\small,
  fontupper=\small,
}

\newcommand{\techtable}[3]{%
\begin{center}
\renewcommand{\arraystretch}{1.3}
\begin{tabularx}{\linewidth}{l X}
\toprule
\textbf{#2} & \textbf{#3} \\
\midrule
#1
\bottomrule
\end{tabularx}
\end{center}
}

\newcommand{\docfigure}[3]{%
  \begin{figure}[h]
    \centering
    \includegraphics[width=#2\linewidth]{#1}
    \caption{\small #3}
  \end{figure}
}

\begin{document}
\pagenumbering{gobble}

% ======================= TITLE PAGE =======================
\begin{titlepage}
\centering
\vspace*{0.5cm}

\includegraphics[width=0.42\linewidth]{aws-iam-monitor-full-logo.png}\\[1cm]

\rule{0.85\linewidth}{0.8pt}\\[0.7cm]

{\Large \color{textgray} Full Technical Documentation}\\[0.5cm]
{\Huge \bfseries \color{primaryblue} aws-iam-monitor}\\[0.4cm]
{\Large \color{primaryblue} Real-time AWS IAM Auditing, Security Monitoring and Automated Threat Detection Platform}\\[0.7cm]

\rule{0.85\linewidth}{0.8pt}\\[1cm]

\begin{center}
\begin{tabularx}{0.82\linewidth}{@{}l X@{}}
\textbf{Version:} & 1.2.0 \\[6pt]
\textbf{Author:} & ZouariOmar \\[6pt]
\textbf{License:} & Apache-2.0 \\[6pt]
\textbf{Repository:} & \url{https://github.com/ZouariOmar/aws-iam-monitor} \\[6pt]
\textbf{Contact:} & \url{zouariomar20@gmail.com} \\
\end{tabularx}
\end{center}

\vfill

{\small \color{textgray} This document consolidates the project's reference documentation --
overview, architecture, installation, usage and security -- into a single,
self-contained technical reference for the \texttt{aws-iam-monitor} v1.2.0
release.}

\end{titlepage}

\pagenumbering{roman}
\tableofcontents
\newpage
\pagenumbering{arabic}

% ======================= CHAPTER: OVERVIEW =======================
\chapter{Overview}

\section{Introduction}

\textbf{aws-iam-monitor} is a production-ready, serverless AWS security
solution designed to provide continuous, real-time visibility into Identity
and Access Management (IAM) changes across AWS accounts and AWS
Organizations.

By capturing management events from \textbf{AWS CloudTrail}, filtering
events with \textbf{Amazon EventBridge}, and processing security payloads
with \textbf{AWS Lambda}, the system detects unauthorized modifications,
privilege escalation risks, new credential creation, and policy tampering
within seconds. High-risk events immediately trigger \textbf{Amazon SNS}
notifications (with optional automated email subscriptions) while
maintaining structured audit logs in \textbf{Amazon S3} and custom
performance metrics in \textbf{Amazon CloudWatch}.

\section{Key Features}

\begin{itemize}[leftmargin=1.4em]
  \item \textbf{Real-time Event Detection}: Instant detection of sensitive IAM operations (e.g. \texttt{CreateUser}, \texttt{DeletePolicy}, \texttt{CreateAccessKey}, \texttt{DeactivateMFADevice}, \texttt{AttachUserPolicy}).
  \item \textbf{IP Whitelisting \& Enforcement}: Automated generation of \texttt{IPWhitelistPolicy} and real-time source IP validation within the monitoring Lambda.
  \item \textbf{Risk Classification}: Automated risk grading (\texttt{CRITICAL}, \texttt{HIGH}, \texttt{MEDIUM}, \texttt{LOW}) for every IAM API call.
  \item \textbf{Dedicated SNS Management Component}: Dedicated \texttt{sns} module with automated email alert subscriptions via \texttt{.env} or CLI parameters.
  \item \textbf{Audit Trail Archiving}: Every IAM event is saved as a structured JSON record in Amazon S3 grouped by date (\texttt{YYYY/MM/DD/}).
  \item \textbf{CloudWatch Metrics}: Emits custom \texttt{AWSIAMMonitor} metrics for CloudWatch dashboards and alarm creation.
  \item \textbf{Idempotent Infrastructure}: Both the Bash CLI (\texttt{awsctl}) and the Terraform implementation create, update, and delete resources safely without duplication.
  \item \textbf{Two Equivalent Implementations}: An imperative Bash/AWS-CLI CLI and a declarative Terraform implementation, pick whichever fits your workflow.
\end{itemize}

\section{Repository Structure}

\begin{lstlisting}
aws-iam-monitor/
|-- Makefile
|-- README.md
|-- ARCHITECTURE.md
|-- INSTALL.md
|-- USAGE.md
|-- SECURITY.md
|-- SUPPORT.md
|-- project/
|   |-- bash/                       # Bash / AWS CLI implementation
|   |   |-- .env.example            # Environment configuration template
|   |   |-- awsctl                  # Primary orchestration CLI binary
|   |   |-- lib/                    # Shared helper library
|   |   |-- iam/                    # IAM module
|   |   |-- cloud-trail/            # CloudTrail module
|   |   |-- sns/                    # SNS module
|   |   |-- event-bridge/           # EventBridge module
|   |   |-- lambda/                 # Lambda module
|   |   |-- Makefile                # Bash implementation Makefile
|   |   `-- README.md               # Bash implementation reference
|   `-- terraform/                  # Terraform implementation
|       |-- main.tf                 # Module wiring
|       |-- variables.tf            # Input variables
|       |-- outputs.tf              # Output values
|       |-- providers.tf            # AWS provider configuration
|       |-- versions.tf             # Terraform & provider version constraints
|       |-- locals.tf               # Computed values
|       |-- terraform.tfvars.example
|       |-- modules/                # iam, cloudtrail, sns, lambda,
|       |                           # eventbridge, s3, cloudwatch
|       |-- Makefile                # Terraform implementation Makefile
|       `-- README.md               # Terraform implementation reference
|-- docapi/                         # Sphinx documentation site source
`-- res/                            # Documentation assets & logos
\end{lstlisting}

\section{Documentation Map}

\techtable{
README.md & Project overview, key features, quickstart. \\
ARCHITECTURE.md & Architecture diagram, AWS services table, deployment ordering. \\
INSTALL.md & Prerequisites, requirements, and setup for both implementations. \\
USAGE.md & Full Bash (\texttt{awsctl}) and Terraform CLI/command reference, deployment examples. \\
SECURITY.md & Security policy, vulnerability reporting, security best practices. \\
project/bash/README.md & Bash implementation deep reference. \\
project/terraform/README.md & Terraform implementation deep reference, including differences from Bash. \\
docapi (Sphinx site) & Full documentation site source, built with \texttt{make docs}. \\
}{Document}{Covers}

% ======================= CHAPTER: ARCHITECTURE =======================
\chapter{Architecture}

This architecture is provisioned identically by both implementations; see
Chapter~\ref{chap:usage} to deploy it via either \texttt{project/bash/}
(imperative) or \texttt{project/terraform/} (declarative).

\section{System Data Flow}

The architecture follows an event-driven serverless pattern:

\begin{center}
\begin{tcolorbox}[colback=lightgray, colframe=primaryblue, arc=2pt, boxrule=0.5pt]
\centering
\textbf{IAM Activity Sources} (users, CI/CD pipelines, or a compromised key) \\
$\downarrow$ \\
\textbf{AWS CloudTrail} (management event capture) \\
$\downarrow$ \\
\textbf{Amazon EventBridge} (IAM API event rule filter) \\
$\downarrow$ \\
\textbf{AWS Lambda -- Python 3.11} (event parsing, IP validation, risk classification) \\
$\downarrow$ \\
\textbf{Amazon SNS} (security alerts) \quad + \quad \textbf{Amazon S3} (JSON audit archival) \quad + \quad \textbf{Amazon CloudWatch} (logs, metrics, alarms)
\end{tcolorbox}
\end{center}

\docfigure{aws-iam-monitor-architecture-diagram.png}{0.95}{Full architecture diagram of \texttt{aws-iam-monitor}, identical across both the Bash and Terraform implementations.}

\section{AWS Services \& Responsibilities}

\techtable{
AWS IAM & Customer-managed policies, groups, users, execution roles, and trust policies. \\
AWS CloudTrail & Captures AWS account management events and API calls across the account. \\
Amazon EventBridge & Evaluates CloudTrail events against IAM pattern filters and invokes Lambda. \\
AWS Lambda & Python 3.11 handler: source IP validation, risk classification, S3 audit logging, SNS alert publishing, CloudWatch metric emission. \\
Amazon SNS & Topic, topic policy, and email subscription management for security alerts. \\
Amazon S3 & Server-side encrypted (AES-256), versioned storage for CloudTrail logs and historical IAM audit records. \\
Amazon CloudWatch & Log group retention management, custom metrics (\texttt{AWSIAMMonitor}), and alarms. \\
}{AWS Service}{Purpose}

\section{Deployment Ordering}

The Bash implementation provisions this architecture as an explicit 5-stage
sequence (\texttt{./awsctl --up all}); the Terraform implementation resolves
the same ordering automatically from its module dependency graph
(\texttt{terraform apply}); both converge on the same end state.

\begin{enumerate}[leftmargin=1.6em]
  \item \textbf{IAM}: Customer-managed policies (\texttt{DevPolicy}, \texttt{ProdPolicy}, \texttt{AuditPolicy}), groups, users, and optional execution roles.
  \item \textbf{CloudTrail}: S3 log bucket and CloudTrail management events trail.
  \item \textbf{SNS}: Alert topic (\texttt{iam-alerts}), topic policy, and email subscription.
  \item \textbf{Lambda}: S3 audit history bucket, execution role, packaged \texttt{lambda\_function.py}, and CloudWatch log group.
  \item \textbf{EventBridge}: Rule pattern registration, Lambda target attachment, and invocation permission.
\end{enumerate}

Teardown (\texttt{./awsctl --down all} / \texttt{terraform destroy}) reverses
this order.

\section{Terraform vs. Bash}

Both implementations provision the same architecture described above; they
are independent, equivalent implementations, not a sequential pipeline versus
a functionally different declarative one. Terraform additionally provisions
optional CloudWatch alarms on the metrics Lambda already emits
(\texttt{UnauthorizedIPAccess} and function \texttt{Errors}), which is a
Terraform-only enhancement, not a divergence in required behavior.

\begin{center}
\renewcommand{\arraystretch}{1.35}
{\small
\begin{tabularx}{\linewidth}{X X X}
\toprule
\textbf{Aspect} & \textbf{Bash / AWS CLI} & \textbf{Terraform} \\
\midrule
Approach & Imperative (explicit command sequence) & Declarative (target state described, order resolved automatically) \\
Lambda IAM permissions & Broad scope (\texttt{S3AuditHistory}, \texttt{Resource: "*"} for SNS) & Strictly scoped to exact resources (reinforced least privilege) \\
S3 public access & Not explicitly blocked & Explicitly blocked on both buckets \\
IAM propagation wait & 6-attempt retry loop with a 10s delay & Deterministic single wait (\texttt{time\_sleep}, 10s) \\
Lambda code deployment & Manual SHA-256 hash comparison & Automatic via \texttt{source\_code\_hash} \\
CloudWatch alarms & Absent (metrics emitted but unused) & Optional, enabled by default \\
Reproducibility & High, dependent on script execution order & High, guaranteed by the dependency graph \\
Maintenance & Simple for small, targeted changes & Facilitated by modularity and declarative state for complex changes \\
\bottomrule
\end{tabularx}
}
\end{center}

\begin{notebox}[Note]
See \texttt{project/terraform/README.md} (section ``Differences from the
Bash implementation'') for the complete comparison. Outside these assumed
differences, resource naming, the EventBridge event pattern, the IAM policy
documents (\texttt{DevPolicy}, \texttt{ProdPolicy}, \texttt{AuditPolicy},
\texttt{IPWhitelistPolicy}) and the SNS topic policy remain strictly
identical between both implementations.
\end{notebox}

\section{Cost Estimate}

Both implementations provision the same resources and therefore share the same
cost profile. The figures below are modeled with the AWS Pricing Calculator
(\texttt{calculator.aws}), Region US East (N.\ Virginia), at standard on-demand
rates with the AWS Free Tier \emph{not} applied, for a baseline of roughly
50{,}000 IAM management events per month.

\begin{center}
\renewcommand{\arraystretch}{1.3}
{\small
\begin{tabularx}{\linewidth}{l X r}
\toprule
\textbf{AWS service} & \textbf{Basis} & \textbf{Est./mo} \\
\midrule
AWS CloudTrail & 50k IAM management events @ \$2.00 / 100k & \$1.00 \\
Amazon EventBridge & AWS service events on the default bus (no charge) & \$0.00 \\
AWS Lambda & 50k requests, 128\,MB, \textasciitilde{}400\,ms (\textasciitilde{}2{,}500 GB-s) & \$0.05 \\
Amazon SNS & 50k publishes + \textasciitilde{}500 email deliveries & \$0.04 \\
Amazon S3 & \textasciitilde{}10\,GB storage, \textasciitilde{}55k PUT/GET & \$0.51 \\
Amazon CloudWatch & 3 custom metrics, 2 alarms, 50k \texttt{PutMetricData}, \textasciitilde{}1\,GB logs & \$2.13 \\
Data transfer & Intra-region + minimal internet egress & \$0.09 \\
\midrule
\textbf{Estimated total} & & \textbf{\$3.82} \\
\bottomrule
\end{tabularx}
}
\end{center}

That is approximately \$3.82 per month (\textasciitilde{}\$46 per year). Cost
scales almost entirely with the IAM event volume: a light workload
(5k events/mo) costs about \$2, a moderate one (50k) about \$4, and a heavy one
(500k) about \$24 per month.

\begin{notebox}[Note]
Amazon CloudWatch dominates the bill (custom metrics plus one
\texttt{PutMetricData} call per event); batching metric publishes reduces it
further. Within the AWS Free Tier most line items fall to \$0 and the effective
total is negligible.
\end{notebox}

% ======================= CHAPTER: INSTALLATION =======================
\chapter{Installation}

\textbf{aws-iam-monitor} ships two independent, equivalent infrastructure
implementations; install prerequisites for whichever you plan to use (or
both). See Chapter~\ref{chap:usage} once installed, for the CLI reference of
both.

\section{Clone Repository}

\begin{lstlisting}
git clone https://github.com/ZouariOmar/aws-iam-monitor.git
cd aws-iam-monitor
\end{lstlisting}

\section{Requirements}

\begin{itemize}[leftmargin=1.4em]
  \item Git
  \item An AWS Account with access to: CloudTrail, EventBridge, Lambda, STS AssumeRole, SNS, S3, CloudWatch
  \item Make (optional, but recommended, wraps both implementations)
\end{itemize}

\subsection{Bash Implementation}

\begin{itemize}[leftmargin=1.4em]
  \item \textbf{Operating System}: Linux / macOS
  \item \textbf{Shell}: Bash v5.0+
  \item \texttt{aws-cli} (v2.0+)
  \item \texttt{jq} (v1.6+)
  \item \texttt{sed} (v4.10+)
  \item \texttt{zip} (v3.0+)
  \item \texttt{openssl} (v3.6+)
  \item \texttt{python3} (v3.11+)
\end{itemize}

\subsection{Terraform Implementation}

\begin{itemize}[leftmargin=1.4em]
  \item Terraform >= 1.5.0
  \item The same AWS account access as above
  \item Network access to the Terraform Registry (\texttt{terraform init} downloads the \texttt{hashicorp/aws}, \texttt{hashicorp/archive}, and \texttt{hashicorp/time} providers automatically, no other local tooling required)
\end{itemize}

\subsection{Documentation Build (optional)}

Python 3.14+ and \texttt{uv} (for the Sphinx \texttt{docapi} site, see
\texttt{make docs}).

\section{Setup}

\subsection{Bash}

\begin{lstlisting}
cd project/bash

# Configure environment variables (optional: SNS email, IP whitelist)
cp .env.example .env
nano .env   # SNS_ALERT_EMAIL=..., ALLOWED_IPS=1.2.3.4/32,5.6.7.8/24

# Ensure CLI executable permissions (or: make -C project/bash chmod)
chmod +x awsctl \
  lib/* \
  iam/src/* \
  cloud-trail/src/* \
  sns/src/* sns/test/* \
  event-bridge/src/* event-bridge/test/* \
  lambda/src/* lambda/test/*

# Verify AWS CLI authentication
aws sts get-caller-identity
\end{lstlisting}

\subsection{Terraform}

\begin{lstlisting}
cd project/terraform
cp terraform.tfvars.example terraform.tfvars
nano terraform.tfvars   # sns_alert_email, allowed_ips, etc.
terraform init
\end{lstlisting}

See \texttt{project/terraform/README.md} for the full reference.

\subsection{Docker}

Each implementation has its own Dockerfile
(\texttt{project/bash/Dockerfile}, \texttt{project/terraform/Dockerfile})
and its own set of Makefile targets, prefixed \texttt{bash-docker-*} and
\texttt{tr-docker-*} respectively:

\begin{lstlisting}
make bash-docker  # build + run + open a shell (Bash implementation)
make tr-docker    # build + run + open a shell (Terraform implementation)
\end{lstlisting}

% ======================= CHAPTER: USAGE =======================
\chapter{Usage}
\label{chap:usage}

\textbf{aws-iam-monitor} ships two independent, equivalent infrastructure
implementations, pick whichever fits your workflow. Do not run both against
the same AWS account simultaneously; they use the same default resource
names and would conflict.

\begin{notebox}[Important]
See Chapter~\ref{chap:usage} above if you haven't set up prerequisites yet
(Installation chapter).
\end{notebox}

\section{Makefile}

The root \texttt{Makefile} delegates Bash- and Terraform-specific commands
to their own sub-Makefiles (\texttt{project/bash/Makefile},
\texttt{project/terraform/Makefile}), which remain independently usable via
\texttt{make -C project/bash <target>} / \texttt{make -C project/terraform
<target>}. Run \texttt{make help} at any level for the full,
self-documenting list.

\subsection{Main}

\techtable{
\texttt{make init} & Init the project (\texttt{uv venv}). \\
\texttt{make sync} & Sync dependencies (\texttt{uv sync}). \\
\texttt{make run} & Run the Bash implementation (\texttt{ARGS=...}). \\
\texttt{make test} & Run the Bash implementation's tests. \\
\texttt{make clean} & Clean the project. \\
\texttt{make changelog} & Generate \texttt{CHANGELOG.md}. \\
\texttt{make version-up} & Bump project version and create a git tag. \\
\texttt{make gh-issue} & Create a GitHub issue from a template (\texttt{TITLE=... FILE=...}). \\
}{Target}{Description}

\subsection{Bash Implementation}

\techtable{
\texttt{make bash-up TARGET=...} & Provision infrastructure (\texttt{TARGET} default: \texttt{all}). \\
\texttt{make bash-down TARGET=...} & Tear down infrastructure. \\
\texttt{make bash-status TARGET=...} & Show infrastructure status. \\
\texttt{make bash-test TARGET=...} & Run the Bash test suite. \\
}{Target}{Description}

\subsection{Terraform Implementation}

\techtable{
\texttt{make tf-init} & \texttt{terraform init}. \\
\texttt{make tf-fmt} & \texttt{terraform fmt}. \\
\texttt{make tf-validate} & \texttt{terraform validate}. \\
\texttt{make tf-plan} & \texttt{terraform plan}. \\
\texttt{make tf-apply} & \texttt{terraform apply}. \\
\texttt{make tf-destroy} & \texttt{terraform destroy}. \\
\texttt{make tf-output} & \texttt{terraform output}. \\
}{Target}{Description}

\subsection{Documentation}

\techtable{
\texttt{make docs} & Build the Sphinx docs to HTML. \\
\texttt{make docs-serve} & Serve the built docs locally. \\
\texttt{make docs-preview} & Build and preview docs (Vercel). \\
\texttt{make docs-deploy} & Build and deploy docs (Vercel). \\
}{Target}{Description}

\subsection{Docker}

\techtable{
\texttt{make bash-docker} & Build + run + open a shell (Bash). \\
\texttt{make bash-docker-build} & Build the Bash Docker image. \\
\texttt{make bash-docker-bash} & Open a shell inside the Bash container. \\
\texttt{make bash-docker-run} & Run \texttt{awsctl} inside the Bash container. \\
\texttt{make bash-docker-stop} & Stop the running Bash container. \\
\texttt{make bash-docker-clean} & Remove the Bash Docker image. \\
\texttt{make tr-docker} & Build + run + open a shell (Terraform). \\
\texttt{make tr-docker-build} & Build the Terraform Docker image. \\
\texttt{make tr-docker-bash} & Open a shell inside the Terraform container. \\
\texttt{make tr-docker-run} & Run \texttt{terraform} inside the Terraform container. \\
\texttt{make tr-docker-stop} & Stop the running Terraform container. \\
\texttt{make tr-docker-clean} & Remove the Terraform Docker image. \\
}{Target}{Description}

\section{Bash -- \texttt{awsctl}}

The primary orchestration tool is \texttt{./awsctl}, run from
\texttt{project/bash/}. It provides a unified command line interface for
provisioning, managing, testing, and destroying infrastructure components.

\subsection{Command Syntax}

\begin{lstlisting}
cd project/bash
./awsctl <ACTION> [TARGET] [OPTIONS]
\end{lstlisting}

\subsection{Actions}

\begin{itemize}[leftmargin=1.4em]
  \item \texttt{--up} : Create/configure resources.
  \item \texttt{--down} : Tear down/delete resources.
  \item \texttt{--status} : Check status of resources.
  \item \texttt{--test} : Execute test suites.
\end{itemize}

\subsection{Targets}

\begin{itemize}[leftmargin=1.4em]
  \item \texttt{all} \textit{(default)} : Complete monitoring infrastructure.
  \item \texttt{iam} : IAM policies, groups, users, and roles.
  \item \texttt{cloud-trail} : CloudTrail trail and log S3 bucket.
  \item \texttt{sns} : SNS security alert topic and email subscriptions.
  \item \texttt{lambda} : Lambda function and audit S3 bucket.
  \item \texttt{event-bridge} : EventBridge rule and target registration.
\end{itemize}

\subsection{Options}

\begin{itemize}[leftmargin=1.4em]
  \item \texttt{-r}, \texttt{--role} : Include optional IAM roles during IAM setup.
  \item \texttt{-v}, \texttt{--verbose} : Display detailed AWS CLI execution logs.
  \item \texttt{-h}, \texttt{--help} : Show help message.
\end{itemize}

\subsection{Deployment Examples}

\begin{lstlisting}
./awsctl --up all          # Provision the complete monitoring pipeline
\end{lstlisting}

\textit{This will automatically:}

\begin{enumerate}[leftmargin=1.6em]
  \item Create customer-managed IAM policies (\texttt{DevPolicy}, \texttt{ProdPolicy}, \texttt{AuditPolicy}), groups (\texttt{DevGroup}, \texttt{ProdGroup}, \texttt{TestGroup}), and users (\texttt{DevAdmin}, \texttt{ProdAdmin}, etc.).
  \item Provision an S3 log bucket and create a CloudTrail management events trail.
  \item Provision the SNS alert topic (\texttt{iam-alerts}) and subscribe \texttt{SNS\_ALERT\_EMAIL} if configured in \texttt{.env}.
  \item Provision an S3 audit history bucket and deploy the monitoring Lambda function with its execution role.
  \item Create an EventBridge rule for IAM events, attach the Lambda function as a target, and grant invocation permissions.
\end{enumerate}

\begin{lstlisting}
./awsctl --up sns          # Provision only the SNS alert component
./awsctl --status all      # Check status of every component
./awsctl --test all        # Run automated test suites
./awsctl --down all        # Tear down the complete infrastructure
\end{lstlisting}

\begin{notebox}[Note]
For prerequisites, setup, module structure, and a detailed comparison
against the Terraform implementation, see \texttt{project/bash/README.md}.
\end{notebox}

\section{Terraform}

The equivalent declarative implementation lives under
\texttt{project/terraform/}.

\subsection{Command Syntax}

\begin{lstlisting}
cd project/terraform
terraform init
terraform plan
terraform apply
terraform destroy
\end{lstlisting}

A plain \texttt{terraform apply} with default variables provisions the same
architecture as \texttt{./awsctl --up all}, including the demo IAM sandbox
(\texttt{enable\_iam\_sandbox = true} by default). Configuration is entirely
variable-driven via \texttt{terraform.tfvars} (copy from
\texttt{terraform.tfvars.example}) -- Terraform does not read the Bash
implementation's \texttt{.env} file.

\subsection{Deployment Examples}

\begin{lstlisting}
# First-time provisioning
cp terraform.tfvars.example terraform.tfvars
terraform init
terraform apply

# Updating infrastructure, edit terraform.tfvars, then:
terraform plan
terraform apply

# Teardown
terraform destroy
\end{lstlisting}

Changing the shared Lambda source
(\texttt{project/bash/lambda/src/lambda\_function.py}) is detected
automatically via \texttt{source\_code\_hash} and redeployed on the next
\texttt{apply}, no manual packaging step.

\begin{notebox}[Note]
For the full variable reference, state security guidance, and a detailed
comparison against the Bash implementation, see
\texttt{project/terraform/README.md}.
\end{notebox}

% ======================= CHAPTER: SECURITY =======================
\chapter{Security}

\section{Supported Versions}

The following versions of \textbf{aws-iam-monitor} receive security updates.

\begin{center}
\renewcommand{\arraystretch}{1.3}
\begin{tabularx}{0.5\linewidth}{X c}
\toprule
\textbf{Version} & \textbf{Supported} \\
\midrule
1.x.x & yes \\
0.x.x & no \\
\bottomrule
\end{tabularx}
\end{center}

\section{Reporting a Vulnerability}

If you discover a security vulnerability in \textbf{aws-iam-monitor}, please
report it responsibly.

\subsection{How to Report}

Send an email to \textbf{zouariomar20@gmail.com} containing:

\begin{itemize}[leftmargin=1.4em]
  \item A description of the vulnerability
  \item Steps to reproduce
  \item The affected version
  \item Any relevant logs or screenshots
  \item A proof of concept (if available)
\end{itemize}

Please \textbf{do not open a public GitHub issue} for security
vulnerabilities.

\subsection{Response Process}

We aim to:

\begin{itemize}[leftmargin=1.4em]
  \item Acknowledge your report within \textbf{48 hours}
  \item Investigate and validate the issue
  \item Keep you informed throughout the remediation process
  \item Release a security fix as quickly as practical
  \item Credit you for the disclosure (if desired)
\end{itemize}

\section{Security Practices}

The project follows security best practices for AWS environments.

\subsection{Least Privilege}

\begin{itemize}[leftmargin=1.4em]
  \item IAM policies grant only the permissions required.
  \item Cross-account access should use read-only roles whenever possible.
  \item AWS STS AssumeRole is recommended for organization-wide deployments.
\end{itemize}

\subsection{IP Whitelisting}

\begin{itemize}[leftmargin=1.4em]
  \item The system supports automated \texttt{IPWhitelistPolicy} generation to restrict IAM actions to trusted networks.
  \item Continuous monitoring validates the source IP of every IAM event against a pre-configured \texttt{ALLOWED\_IPS} list.
  \item Unauthorized access attempts from blocked IPs trigger high-priority alerts and custom CloudWatch metrics.
\end{itemize}

\subsection{Idempotent, Safe Operations}

\begin{itemize}[leftmargin=1.4em]
  \item Both implementations check resource state before creating, updating, or deleting it, to avoid duplicate resources and unintended service disruptions.
  \item The Bash implementation's \texttt{awsctl} and its \texttt{*\_ctl} sub-scripts verify existence/hashes before acting; the Terraform implementation relies on Terraform's own state-diffing (\texttt{terraform plan}) for the same guarantee, a second \texttt{terraform plan} after \texttt{apply} should show no changes.
\end{itemize}

\subsection{Secure AWS Integrations}

aws-iam-monitor relies on AWS managed services including AWS CloudTrail,
Amazon EventBridge, AWS Lambda, Amazon SNS, Amazon S3, Amazon CloudWatch, and
AWS STS. These services should be configured according to AWS security best
practices.

\subsection{Data Protection}

\begin{itemize}[leftmargin=1.4em]
  \item Communication with AWS APIs uses HTTPS/TLS.
  \item Audit reports stored in Amazon S3 should use encryption at rest.
  \item Sensitive configuration should never be committed to the repository.
\end{itemize}

\subsection{Secrets Management}

Never commit: AWS Access Keys, Secret Access Keys, Session Tokens,
\texttt{.env} files, private certificates, or customer data.

Use one of the following instead: IAM Roles, AWS Secrets Manager, AWS
Systems Manager Parameter Store, or environment variables.

\subsection{Logging}

aws-iam-monitor records IAM activity for auditing purposes. Logs should
exclude secrets and credentials, be retained according to organizational
policies, and be monitored for suspicious activity.

\subsection{Dependency Management}

Dependencies should be kept up to date. Recommended tools include
Dependabot, \texttt{pip-audit}, \texttt{uv}, and GitHub Dependabot Alerts.

\section{Security Contributions}

Contributors can help improve security by following secure coding
practices, validating all external input, avoiding hard-coded credentials,
using least-privilege IAM permissions, reviewing dependency updates, and
testing security-sensitive changes before submitting pull requests.

\section{AWS Security Recommendations}

For production deployments, we recommend:

\begin{itemize}[leftmargin=1.4em]
  \item Enable AWS CloudTrail in all regions.
  \item Enable S3 bucket encryption.
  \item Enable CloudWatch log retention.
  \item Use SNS encryption where appropriate.
  \item Enable MFA for privileged AWS accounts.
  \item Rotate IAM credentials regularly.
  \item Use dedicated monitoring accounts within AWS Organizations.
\end{itemize}

\section{Resources}

\begin{itemize}[leftmargin=1.4em]
  \item AWS Well-Architected Framework -- Security Pillar
  \item AWS IAM Best Practices
  \item AWS Security Best Practices
  \item OWASP Secure Coding Practices Checklist
  \item CWE Top 25 Most Dangerous Software Weaknesses
\end{itemize}

% ======================= CHAPTER: RESULTS =======================
\chapter{Results \& Illustrations}

The following screenshots, captured on the AWS test account used during
development, illustrate the system's effective behavior once deployed: an
IAM event is captured by CloudTrail, filtered by the EventBridge rule,
analyzed and classified by the Lambda function, archived to the audit
bucket, and -- for high-risk events -- notified via the SNS topic
\texttt{iam-alerts}.

\section{SNS Subscription Confirmation}

\docfigure{mail/aws-notification-subscription-confirmation.png}{0.62}{Confirmation e-mail received when subscribing to the \texttt{iam-alerts} SNS topic.}

\section{High-Risk Alerts}

\docfigure{mail/hight-risk-aws-iam-security-alert-create-user.png}{0.75}{Alert notification for an event classified as high risk: IAM user creation (\texttt{CreateUser}).}

\docfigure{mail/hight-risk-aws-iam-security-alert-delete-user.png}{0.75}{Alert notification for an event classified as high risk: IAM user deletion (\texttt{DeleteUser}).}

\section{Unauthorized IP Access Detection}

\docfigure{mail/security-alert-unauthorized-ip-access-attempt-detected.png}{0.75}{Alert triggered when an IAM operation originates from a source IP absent from the configured whitelist.}

% ======================= CHAPTER: CHANGELOG =======================
\chapter{Release Notes}

\section{v1.2.0 -- Current Release}

\begin{itemize}[leftmargin=1.4em]
  \item \textbf{Features}: Makefile GitHub issue target (\#13); Terraform AWS infrastructure implementation.
  \item \textbf{Other}: Support for Bash submodule in the root build.
  \item \textbf{Refactor}: Bash tooling moved into its own subdirectory.
  \item \textbf{Documentation}: v1.1.0 changelog; updated architecture diagram; Terraform documentation added to \texttt{docapi}; project documentation updates.
\end{itemize}

\section{v1.1.0}

\begin{itemize}[leftmargin=1.4em]
  \item \textbf{Features}: Lambda source IP whitelist enforcement; IAM whitelist policy provisioning; improved Docker workflow and Makefile.
  \item \textbf{Bug Fixes}: Makefile version handling; SNS subscription state handling; CloudTrail S3 status checks region fix.
  \item \textbf{Refactor}: Modernized Lambda timestamps and security alerts.
  \item \textbf{Documentation}: Security documentation update.
\end{itemize}

\section{v1.0.0}

\begin{itemize}[leftmargin=1.4em]
  \item \textbf{Features}: EventBridge \texttt{IAM CloudTrail} event pattern policy and rule management controller -- first stable release of the complete monitoring pipeline (IAM, CloudTrail, EventBridge, Lambda, SNS, S3, CloudWatch).
\end{itemize}

\begin{notebox}[Note]
See \texttt{CHANGELOG.md} at the repository root for the complete,
auto-generated release history (generated with \texttt{make changelog}).
\end{notebox}

% ======================= BACK MATTER =======================
\chapter*{License \& Contact}
\addcontentsline{toc}{chapter}{License \& Contact}
\markboth{License \& Contact}{License \& Contact}

\section*{License}
\addcontentsline{toc}{section}{License}

This project is licensed under the \textbf{Apache-2.0 License}. See the
\texttt{LICENSE} file at the repository root for details.

\section*{Contact \& Support}
\addcontentsline{toc}{section}{Contact \& Support}

\begin{itemize}[leftmargin=1.4em]
  \item \textbf{Author}: ZouariOmar (\url{https://github.com/ZouariOmar})
  \item \textbf{Email}: \url{zouariomar20@gmail.com}
  \item \textbf{LinkedIn}: \url{https://www.linkedin.com/in/zouari-omar}
  \item \textbf{Issues}: \url{https://github.com/ZouariOmar/aws-iam-monitor/issues}
  \item \textbf{Support}: See \texttt{SUPPORT.md}
\end{itemize}

\end{document}
